\documentclass{beamer}


\usepackage{tikz}
\usepackage[linguistics]{forest}
\usepackage[many]{tcolorbox}
\usepackage{comment}
\usepackage{multicol}
\usepackage{listings}
\usepackage{prooftree}

\newif\ifframeinlbf
\frameinlbftrue
\makeatletter
\newcommand\listofframes{\@starttoc{lbf}}
\makeatother

\addtobeamertemplate{frametitle}{}{%
    \ifframeinlbf%
         \addtocontents{lbf}{%
            \protect\makebox[2em][l]{%
                \protect\usebeamercolor[fg]{structure}\insertframenumber\hfill%
            }%
            \protect\hyperlink{page.\insertpagenumber}\insertframetitle\par%
        }%
   \else\fi%
}

\AtBeginSection[]{
  \begin{frame}
  \vfill
  \centering
  \begin{beamercolorbox}[sep=8pt,center,shadow=true,rounded=true]{title}
    \usebeamerfont{title}\insertsectionhead\par%
  \end{beamercolorbox}
  \vfill
  \end{frame}
}

\newcommand{\myalert}[2]{%
    \tikz\node[coordinate](t#1){};\alert<#1>{#2}
}


\newcommand{\myexplain}[3]
{\begin{tcolorbox}
    [   tikznode boxed title,
        enhanced,
        arc=0mm,
        interior style={white},
        attach boxed title to top center= {yshift=-\tcboxedtitleheight/2},
        fonttitle=\bfseries,
        colbacktitle=white,coltitle=red,
        boxed title style={size=small,colframe=red,boxrule=0pt},
        title={\tikz\node[coordinate](s#1) {};#2}]
#3
\end{tcolorbox}
\begin{tikzpicture}[overlay]
\path<#1>[->, thick] (s#1) edge [bend left] (t#1);
\end{tikzpicture}
}

\newcommand{\explain}[2]
{\begin{tcolorbox}
    [   tikznode boxed title,
        enhanced,
        arc=0mm,
        interior style={white},
        attach boxed title to top center= {yshift=-\tcboxedtitleheight/2},
        fonttitle=\bfseries,
        colbacktitle=white,coltitle=red,
        boxed title style={size=small,colframe=red,boxrule=0pt},
        title={#1}]
#2
\end{tcolorbox}
}

\def\point{\begin{tikzpicture}
\fill[orange] (0,0) circle (2pt);
\end{tikzpicture}}

\setbeamertemplate{navigation symbols}{} %Remove navigation bar

\tikzstyle{every picture}+=[remember picture]

\usepackage{color}
\newcommand{\vet}[1]{\foreach \num in {#1}{\el{\num}}}
\newcommand{\el}[1]{\fbox{\parbox[c][2ex][c]{1.1em}{#1}}}
\newcommand{\id}[1]{\parbox[r][2ex][b]{1.7em}{\raggedleft #1}}
\newcommand{\bl}{\color{blue}}
\newcommand{\gr}{\color{gray}}
\newcommand{\re}{\color{red}}

\title{Regular Expressions}
\author{Lunjin Lu}
\date{}


\begin{document}
\begin{comment}
\end{comment}
\frameinlbffalse
\begin{frame}[fragile,t,allowframebreaks]{List of Frames}
    \listofframes
\end{frame}
\frameinlbftrue


\frame{ \titlepage
}

\begin{frame}{Regular Expressions}
Regular Expressions are patterns for strings. A regular expression define a langauge that 
consists of those strings that conform to the pattern. For example
\begin{itemize}
  \item a 
  \item b*
  \item ab*
  \item ab*+ba*
\end{itemize}
\end{frame}

\begin{frame}{Grammar of Regular Expressions }
The set of regular expressions RegEx is inductively defined as follows. 
\begin{itemize}
  \item $\Lambda\in RegEx$
  \item $\emptyset\in RegEx$
  \item $\ldots$
\end{itemize}
\end{frame}

\begin{frame}{Example Regular Expressions}

\end{frame}

\begin{frame}{Language Concatenation}

\end{frame}

\begin{frame}{Language of a Regular Expression} 
The language of a regular expresison $r$ is denoted as $L(r)$ which is defined as follows. 
\begin{itemize}
  \item \ldots
\end{itemize}

A language $L$ is said to be reglar iff there is a regular expression such that $L=L(r)$.
\end{frame}

\begin{frame}{Equivalence between Regular Expressions}
Two regular expressions are equivalent if their langauges are the same. 
\[ r_1=r_2 \Longleftrightarrow L(r_1)=L(r_2)
\]
\end{frame}

\end{document}